\scanpage{154}
\mspage{12}
% archive sheet number: 154
de $\mathrm{Top}(T)^\wedge$, i.e. les faisceaux d'ens. \inserted{$F$} 0-connexes sur $T$, sont ceux dont la ``fibre générique'' a exactement un point \unsure{...|ou.})
\inserted{Ainsi $\mathscr T=\mathrm{Top}(T)$ est loc. irr. ssi $T$ loc. connexe, et les comp. connexes de $T$ sont \unsure{irr.} irréductibles, i.e. ssi $T=\coprod_i T_i$, les $T_i$ irréductibles.}

Si $\mathscr T$ est de la forme $\mathrm{Top}(X)$, $X$ une catégorie, \unsure{où|on} sait que $\mathscr T$ est loc. connexe (et \unsure{$=^t$} localement $\infty$-connexe).
On voit \deleted{donc} que $\mathscr T$ est irréductible ssi $\mathrm{Esp}(X)$ (l'espace topologique associé, dont l'ens. sous-jacent est $\mathrm{Ob}\,X$) est irréductible\inserted{*}.
% slanted
\margin{* Car les ouverts de $\mathrm{Top}(X)$, ceux de $\mathrm{Esp}(X)$ et ceux de $X$, \unsure{s'identifient}. (\deleted{\unsure{NB}} Cela signifie (ssi) que pour $x,y$ dans $X$, $\exists z$ \unsure{avec}
\diagram{\begin{tikzcd}[ampersand replacement=\&,row sep=0pt,column sep=small]
 \& x\\ z\arrow[ur]\arrow[dr] \& {}\\ \& y
\end{tikzcd}}
\unsure{;} i.e. $\mathrm{Ord}(X)$ est filtrant à \unsure{g.}}
On \deleted{\unsure{définira}} dit alors que $X$ est irréductible. \deleted{\unsure{ssi}} Comme les composantes connexes de $X$ et celles de $\mathrm{Esp}(X)$ \inserted{ou $\mathrm{Top}(X)=\mathscr T$} se correspondent, on voit donc que \deleted{\unsure{les}} les conditions suivantes sont équivalentes :
\textcircled{a} \inserted{le topos} $\mathrm{Top}(X)$ est loc. irr. \quad b) $\mathrm{Esp}(X)$ est loc. irréductible \quad c) les composantes connexes de \deleted{\unsure{objets}} $X$ sont irréductibles\unsure{;} \inserted{\textcircled{d} les $X_{/x}$ irréductibles,}
On dit alors que $X$ est \deleted{loc.} irréductible.
% slanted
\margin{Faux, on a
\diagram{\begin{tikzcd}[ampersand replacement=\&,row sep=small]
 (a\Longleftrightarrow d)\arrow[d,Rightarrow]\\
 (b\Longleftrightarrow c)
\end{tikzcd}}}
% slanted
\margin{$X$ irréductible n'implique \uline{pas} que $X_{/x}$ le soient, exemple
\diagram{\begin{tikzcd}[ampersand replacement=\&,column sep=small]
 X={} \& 0\arrow[r,bend left,"\alpha"]\arrow[r,bend right,"\beta"'] \& 1 \& (\alpha\ne\beta),
\end{tikzcd}}
on a
\diagram{\begin{tikzcd}[ampersand replacement=\&,row sep=0pt,column sep=small]
 \& 0'\arrow[dr]\\ X_{/1}=\deleted{\unsure{0}} \& \& 1\\
 \&0''\arrow[ur]
\end{tikzcd}}
\unsure{qui} n'est \uline{pas} irréductible, \deleted{\unsure{et n'est pas localement irréductible}} ni \unsure{localt irrd.|localt itou} puisqu'il est connexe. Pourtant il est étale sur $X$.}
\deleted{Il \unsure{reviendra}} Cela implique donc que pour \deleted{tout objet \unsure{0-connexe}} $X'$ étale sur $X$, $X'$ 0-connexe, $X'$ est irréductible.
\inserted{\uline{Dire conditions équivalentes} \textcircled{a} $\mathrm{Top}(X)$ est loc. irr.
\textcircled{b} les $X_{/x}$ irréductibles \textcircled{c} \deleted{tout} objet 0-connexe de $X^\wedge$ est irréductible (i.e. \unsure{tt} \unsure{s.-obj.} $\ne\varnothing$ \unsure{etc}),
\deleted{\textcircled{d} les \unsure{commas} de $X$ \unsure{sont irréductibles}} \textcircled{e} conditions de \deleted{\unsure{SGA}} $\mathrm{SGA}\,4$ \unsure{I 2|I § 2}}
Si \deleted{$X$ est \unsure{Cat.} \unsure{vide}} \inserted{$X\ne\varnothing$ et} $X$ irréductible signifie \deleted{\unsure{ssi}} que $\forall\,x,y\in\mathrm{Ob}\,X$, existe $s$ \deleted{\unsure{et deux}}
\diagram{\begin{tikzcd}[ampersand replacement=\&,row sep=0pt,column sep=small]
 \& x\\ s\arrow[ur]\arrow[dr] \& {}\\ \& y
\end{tikzcd}}
. (Cela signifie que $X$ est \uline{fbt. filtrante}* à gauche, faiblement cofiltrante).
$X$ loc. \deleted{\unsure{filtr}} irréd signifie que pour $x,y$ \deleted{\unsure{qq}} \uline{appartenant à une \unsure{même} comp. connexe}, existe $s$ avec
\diagram{\begin{tikzcd}[ampersand replacement=\&,row sep=0pt,column sep=small]
 \& x\\ s\arrow[ur]\arrow[dr] \& {}\\ \& y,
\end{tikzcd}}
% slanted
\margin{* Il vaut mieux l'appeler \uline{pseudo-filtrante}, et réserver ``loc. \unsure{filt.}'' (= fbt. filtrante $\unsure{+}$ connexe).}
